\begin{abcliste}
\abc
\begin{align}
 E_\mathrm{pot} &= \EpF = \mS\cdot\g\cdot\hS = \Ep \approx \result{\EpP{3}{2}}\\
 E_\mathrm{kin} &= \EkF = \frac{1}{2}\cdot\mS\cdot\left(\vS\right)^2 = \Ek \approx \result{\EkP{3}{2}}
\end{align}

\abc
Die Energiebilanz mit Reibung lautet $m\,g\,h = \frac{1}{2}\,m\,v^2 + E_\mathrm{th}$:
\begin{align}
 E_\mathrm{th} &= \EthF = \Ep - \Ek = \Eth \approx \result{\EthP{3}{2}}
\end{align}
Das sind
\begin{align}
 \frac{E_\mathrm{th}}{E_\mathrm{pot}} &= \frac{\Eth}{\Ep} \approx \result{\qthP{0}{2}\,\%}
\end{align}
der Lageenergie. Ohne Reibung wäre die ganze Lageenergie zu kinetischer Energie geworden:
\begin{align}
 v_0 &= \voF = \sqrt{2\cdot\g\cdot\hS} = \vo \approx \result{\voP{0}{2}}
\end{align}

\abc
Die thermische Energie steckt in den Kufen und im Schnee (und ein wenig in der Luft): Beim Reiben werden sie etwas wärmer. Mit der Zeit gibt sie die Wärme an die Umgebung ab.

\abc
Energie wird weder vernichtet noch verbraucht, sie wird nur umgewandelt. Die Gesamtenergie bleibt gleich: Am Start \EpP{3}{2} Lageenergie, unten \EkP{3}{2} kinetische und \EthP{3}{2} thermische Energie. Eine Kraft braucht auch keine Energie, um zu wirken (eine Tischplatte trägt ein Buch, ohne Energie zu verbrauchen). Die Reibungskraft verrichtet Arbeit gegen die Bewegung; diese Reibungsarbeit ist genau die erzeugte thermische Energie.
\end{abcliste}
